\documentclass[12pt,a4paper]{article}
\usepackage{geometry}
\geometry{left=2.5cm,right=2.5cm,top=2.0cm,bottom=2.5cm}
\usepackage[english]{babel}
\usepackage{amsmath,amsthm}
\usepackage{amsfonts}
\usepackage[longend,ruled,linesnumbered]{algorithm2e}
\usepackage{fancyhdr}
\usepackage{ctex}
\usepackage{array}
\usepackage{listings}
\usepackage{color}
\usepackage{graphicx}

\begin{document}


\title{
{《优化理论与算法》第 {$4$} 次作业
}
}
\date{}

% \author{
% 姓名：\underline{你的名字}~~~~~~
% 学号：\underline{你的学号}~~~~~~}

\maketitle

\begin{itemize}
	\item 作业截止于{\bf 周五（4/3/2026）}，在Canvas系统中提交PDF或照片文件
	\item 作业题目的tex文件可以在Canvas中下载
	\item 鼓励相互讨论，但不要抄袭.
\end{itemize}
%%%%%%%%%%%%%%%%%%%%%%%%%%%%%%%%%%%%%%%%%%%%%%%%%%%%%%%%%%%%%%%%%
%%%%%%%%%%%%%%%%%%%%%%%%%%%%%%%%%%%%%%%%%%%%%%%%%%%%%%%%%%%%%%%%%

\section{阅读与积累}
\paragraph{1.1}
\begin{itemize}
	\item 阅读Bertsimas and Tsitsikils，也就是Introduction to Linear Optimization教材，第3章与第4.1-4.4节相关内容。
\end{itemize}

\section{习题}


\paragraph{2.1}
在利用表格法求解线性规划过程中，假设当前基为 $B=\{3,4,5\}$有如下表格：
\begin{center}
	\begin{tabular}{|c| c c c c c |c|} 
		\hline
		B & $\delta$ &-2 &0 &0 &0 & -10 \\
		\hline
		3 & -1 & $\eta$ &1 &0 &0 & 4 \\
		4 & $\alpha$ & -4 &0 &1 &0 & 1 \\
		5 & $\gamma$ & 3 &0 &0 &1 & $\beta$ \\
		\hline
	\end{tabular}
\end{center}
其中，$\alpha$, $\beta$, $\gamma$, $\delta$, $\eta$为表中未知参数。请给出下面命题成立时$(\alpha, \beta, \gamma, \delta, \eta)$的取值范围
\begin{enumerate}
	\item 当前解是可行解，但不是最优解
	\item 问题是无界问题
	\item 当前解是最优解
\end{enumerate}

\paragraph{2.2}
利用两阶段法求解如下线性规划

\begin{equation*}
\begin{array}{lllll@{}ll}
\text{minimize}  & -3x_1&+x_2&-2x_3 &\\
\text{subject to}& x_1 &+3x_2  &+x_3  &\leq 5 \\
&2x_1 &-x_2 & +x_3  & \geq 2  \\
&4x_1 &+3x_2 &-2x_3   & =5  \\
&x_1 \geq 0, &x_2 \geq 0, & x_3 \geq 0 &  
\end{array}
\end{equation*}

\paragraph{2.3}
利用两阶段法求解如下线性规划

\begin{equation*}
\begin{array}{lllll@{}ll}
\text{maximize}  & x_1&-x_2&+x_3 &\\
\text{subject to}& 2x_1 &-x_2  &-2x_3  &\leq 4 \\
&2x_1 &-3x_2 & -x_3  & \leq -5  \\
&-x_1 &+x_2 &+x_3   & \leq -1  \\
&x_1 \geq 0, &x_2 \geq 0, & x_3 \geq 0 &  
\end{array}
\end{equation*}

\end{document}
